\vskip2mm
\begin{minipage}{0.59\linewidth}
    Pour déterminer la célérité d'une onde ultrasonore de fréquence $N$ dans
    deux milieux différents, on utilise un dispositif constitué d'un émetteur E
    et d'un récepteur R fixés aux extrémités d'un tube.
    E et R sont reliés à un oscilloscope.

    \begin{donnees}
        \begin{itemize}[itemsep=-5pt]
            \item Distance émetteur-récepteur~: $D=\mathrm{ER}=\qty{1}{\m}$~;
            \item $N=\qty{40}{\kHz}$.
        \end{itemize}
    \end{donnees}
\end{minipage}
\hfill
\begin{minipage}{0.40\linewidth}
    \flushright
    \includegraphics[width=\textwidth]{2025_05_09_e1619f2897d39ae202a1g-4(1)}
\end{minipage}


\begin{enumerate}
    \item L'onde ultrasonore est-elle une onde longitudinale ou transversale~?
    \item On remplit le tube par de l'eau. L'oscillogramme ci-contre représente
          le signal émis par E et celui reçu par R.

          Recopier sur votre copie le numéro de la question et écrire la lettre
          correspondante à la proposition vraie.
          \begin{enumerate}
              \item La célérité des ultrasons dans l'eau vaut~:
                    \begin{table}[H]
                        \flushright
                        \setlength{\arrayrulewidth}{1pt}
                        \renewcommand{\arraystretch}{1.5}
                        \begin{tabularx}{.90\textwidth}{|p{3mm}|X|p{3mm}|X|p{3mm}|X|p{3mm}|X|}
                            \hline
                            A & $c=\qty{1520}{\m\per\s}$ &
                            B & $c=\qty{620}{\m\per\s}$ &
                            C & $c=\qty{1667}{\m\per\s}$ &
                            D & $c=\qty{330}{\m\per\s}$ \\
                            \hline
                        \end{tabularx}
                    \end{table}
              \item La longueur d'onde de l'onde ultrasonore vaut~:
                    \begin{table}[H]
                        \flushright
                        \setlength{\arrayrulewidth}{1pt}
                        \renewcommand{\arraystretch}{1.5}
                        \begin{tabularx}{.90\textwidth}{|p{3mm}|X|p{3mm}|X|p{3mm}|X|p{3mm}|X|}
                            \hline
                            A & $\lambda=\qty{25.2}{\mm}$ &
                            B & $\lambda=\qty{30.5}{\mm}$ &
                            C & $\lambda=\qty{37.2}{\mm}$ &
                            D & $\lambda=\qty{41.7}{\mm}$ \\
                            \hline
                        \end{tabularx}
                    \end{table}
          \end{enumerate}

    \item On remplace l'eau par un autre liquide, on constate que le décalage
          horaire entre le signal émis et le signal reçu est
          $\Delta t=\qty{0.9}{\s}$.

          La célérité des ultrasons dans le liquide, a-t-elle augmenté ou
          diminué par rapport à celle dans l'eau~? Justifier.
\end{enumerate}

